\reviewsection{Module 1 Changed How I Interpret; Module 2 Changed What I Must Do}

Module 1 journal argued that the world comes to understand developmental difference through labels, scientific claims, stories, school routines, family knowledge, media, relationships, and the voices institutions choose to treat as credible \citep{garciabautista2026module1}. Cole and Telesford gave me useful diagnostic and educational language while emphasizing autism's heterogeneous presentation \citep{cole2023autism}. Ortega showed how neurological language can become identity, community, and resistance while still becoming reductive when it is allowed to explain the whole person \citep{ortega2013cerebralizing}. The National Education Association connected inclusion to concrete communication, sensory, instructional, and social supports, while \textit{Autism: Insight from Inside} showed how people sharing a diagnosis can have very different lives and support needs \citep{nea2020guide,infobase2016inside}.

Module 2 begins after that interpretation. Once an adult notices a behavior, selects a label, forms a hypothesis, or chooses a support, what happens to the student's power? Does the strategy make communication more dependable, instruction more accessible, and refusal safer? Or does it make the student less visible while leaving the original barrier intact? A plan can reduce a behavior that bothers adults and still reduce the learner's trust, bodily autonomy, cultural identity, self-advocacy, or access to rigorous learning. That distinction is personal for me. As a Latino Dominican, bilingual, neurodivergent, and chronically ill learner and educator, I know how quickly an access barrier can be rewritten as a character flaw. Fatigue can become lack of commitment. Difficulty with an inaccessible process can become irresponsibility. Communication difference can become evidence that the thinking matters less. Positive support cannot require students to hide pain, uncertainty, sensory needs, language differences, or another way of participating to be judged capable.

My answer to this module's question is therefore relational rather than procedural. Support strategies respect dignity, consent, and neurodiversity when the learner has meaningful influence over the goal, communication route, conditions, and definition of success; when adults treat behavior as evidence rather than identity; and when the team is accountable for changing its own interpretation, instruction, and environment.
